\section{Areal--volumetric transition and spectral scale}
\label{sec:ym-areal-volumetric-scale}

The constant determining the gauge semigroup is obtained before the
Hamiltonian is introduced. Its origin combines two independent readers of the same
genealogical state: the return frequency yielding
\(\pi_{\rm HMT}\) and the stationary distribution of the two
areal--volumetric sheets.

\subsection{Dynamics of the two sheets}

Let
\begin{equation}
 F_{\rm av}=\begin{pmatrix}0&1\\1&1\end{pmatrix},
 \qquad e_{\rm ar}=\binom10,
 \qquad e_{\rm vol}=\binom01.
 \label{eq:ym-av-transition-matrix}
\end{equation}
The Fibonacci recurrence gives, by induction,
\begin{equation}
 F_{\rm av}^{n}
 =\begin{pmatrix}F_{n-1}&F_n\\F_n&F_{n+1}\end{pmatrix}.
 \label{eq:ym-av-fibonacci-powers}
\end{equation}
The Perron vector is proportional to \((1,\varphi)\). Since
\(F_{\rm av}\) is symmetric, the Parry measure assigns the two sheets
weights proportional to \((1,\varphi^2)\), and hence
\begin{equation}
 \frac{\mu_{\rm ar}}{\mu_{\rm vol}}=\varphi^{-2},
 \qquad
 \mu_{\rm vol}=\frac{\varphi^2}{1+\varphi^2}.
 \label{eq:ym-av-parry-weights}
\end{equation}
Here \(\mu_{\rm vol}\) is a dimensionless frequency character; it does not
denote a volume of dimension \(L^4\).

The regional reader of \(\pi\) is introduced exactly once through
\begin{equation}
 D_\pi
 =\pi_{\rm HMT}|e_{\rm ar}\rangle\langle e_{\rm ar}|
  +|e_{\rm vol}\rangle\langle e_{\rm vol}|.
 \label{eq:ym-pi-regional-reader}
\end{equation}
The ratio obtained at depth \(n\) is
\begin{equation}
 \Theta_n
 :=\frac{\langle e_{\rm ar},D_\pi F_{\rm av}^{n}e_{\rm ar}\rangle}
 {\langle e_{\rm vol},F_{\rm av}^{n}e_{\rm vol}\rangle}
 =\pi_{\rm HMT}\frac{F_{n-1}}{F_{n+1}}
 \longrightarrow
 r_{\rm av}:=\frac{\pi_{\rm HMT}}{\varphi_{\rm HMT}^{2}}.
 \label{eq:ym-target-free-av-limit}
\end{equation}
The matrix, route, and quotient are fixed before consulting a mass
gap or a gauge-theory parameter. At the native depth \(n=108\),
\begin{equation}
 F_{107}=10284720757613717413913,
 \qquad
 F_{109}=26925748508234281076009,
 \label{eq:ym-fibonacci-107-109}
\end{equation}
and
\begin{equation}
 \left|\Theta_{108}
 -\frac{\pi_{\rm HMT}}{\varphi_{\rm HMT}^{2}}\right|
 <2\cdot10^{-45}.
 \label{eq:ym-av-depth-108-error}
\end{equation}

\subsection{Dimensional transduction}

Let \(\ell_0\) be the length unit of the dimensional reader. The inverse
length scale and the parameter of a slab of thickness \(a_m\) are
\begin{equation}
 \kappa_{\rm av}
 :=\frac{\mu_{\rm vol}}{\ell_0}
 \left(\frac{\pi_{\rm HMT}}{\varphi_{\rm HMT}^{2}}-1\right)>0,
 \qquad
 q_m:=e^{-3a_m\kappa_{\rm av}}.
 \label{eq:ym-dimensional-gap-scale}
\end{equation}
Thus \([\kappa_{\rm av}]=L^{-1}\), and the energy gap is
\begin{equation}
 \Delta_E=3\hbar c\,\kappa_{\rm av};
 \label{eq:ym-energy-gap-units}
\end{equation}
in units \(\hbar=c=1\), this is written as \(\Delta_E=3\kappa_{\rm av}\).
The electronic memory \(\Delta_4\) belongs to a different domain and a different
reader; it does not enter \eqref{eq:ym-dimensional-gap-scale}.

In the semigroups of the following sections, \(t,s,a_m\) are
Euclidean lengths and the generators \(D\) and \(H^{\rm ov}\) have
dimension \(L^{-1}\). The energy Hamiltonian is
\(H_E=\hbar_{\rm int}c_{\rm int}D\). If
\(\tau=t/c_{\rm int}\), then
\[
 e^{-tD}=e^{-\tau H_E/\hbar_{\rm int}},\qquad
 \Delta=3\kappa_{\rm av},\quad
 \Delta_E=\hbar_{\rm int}c_{\rm int}\Delta,\quad
 m_{\rm gap}=\frac{\hbar_{\rm int}\Delta}{c_{\rm int}}.
\]
The same gap is thus expressed, respectively, as inverse
length, energy, and mass in the declared dimensional chart.

Since \(a_m=9a_{m+1}\), the definition immediately gives
\begin{equation}
 q_{m+1}^{9}=q_m.
 \label{eq:ym-nonadic-gap-root}
\end{equation}
This identity is the nonadic root that will transport the same semigroup and
the same spectral scale through every refinement.
